% Living Showcase: Formeln, Math-API und Goal-System (Palette Slot 3).
\StartChapter[ch:math]{Formeln und Zielketten}

\SubsectionBar[sec:math-base]{Formelbox und Math-API}
\ZSFindex{formulabox}
\ZSFindex{Math-API}

\textVorBox{textVorBox bindet Kontext an die folgende Box.}
\begin{formulabox}[Was die Formelbox kann\ZSFTitleTag{formulabox}]
  \textVorBox{textVorBox — innerhalb der Formelbox bindet es an die Formel:}
  \[
    x\in\R,\quad z\in\C,\quad n\in\N,\quad k\in\Z,\quad q\in\Q
  \]
  \formulasep
  \[
    \dd x,\quad \norm{\vect v},\quad \abs{x},\quad \sgn(x),\quad
    \grad f,\quad \divg \vect F,\quad \rot \vect F
  \]
  \formulasep
  \ZSFItemHeading{ZSFItemHeading gliedert innerhalb der Box}
  \[
    \ZSFsumAuto{i=1}{n} i,\quad \ZSFlimAuto{x\to0}\frac{\sin x}{x}
  \]
  \textNachBox{textNachBox — hier die Zeile direkt unter der Formel.}
  \formulanote{formulanote setzt die Anmerkung als eigene Zeile darunter.}
\end{formulabox}
\textNachBox{textNachBox bleibt optisch an der Box gebunden.}

\begin{formulabox}[Hervorhebung, Klammer, Zeilenanmerkung]
  \[
    \ZSFmhlA{a^2}+2\ZSFmhlB{ab}+\ZSFmhlD{b^2}
    =\ZSFmhlC{(a+b)^2}
  \]
  \formulasep
  \[
    \ZSFbraceunder{x^2+2x+1}{ZSFbraceunder}
    =\ZSFbraceover{(x+1)^2}{ZSFbraceover}
  \]
  \formulasep
  \ZSFformulaline{a^2+b^2=c^2}{ZSFformulaline: Anmerkung auf der Zeile}
\end{formulabox}

\SubsectionBar[sec:goal]{Bedingung, Schritt, Ziel}
\ZSFindex{goalbox}

\begin{goalbox}[Zielkette in zwei Schreibweisen\ZSFTitleTag{goalbox}]
  \ZSFDerivationCase*{ZSFDerivationCase* — einzeilig}{\sqrt{x^2}}{x}
  \formulasep
  \ZSFDerivationCase{ZSFDerivationCase}{\sqrt{x^2}}{-x}
  \formulasep
  \ZSFItemHeading{Frei aus den Primitiven gebaut}
  \GoalCondition{GoalCondition}
  \GoalStep{\text{GoalStep}}
  \GoalStep{\text{beliebig oft wiederholbar}}
  \GoalTarget{\text{GoalTarget}}
\end{goalbox}

\begin{statementbox}[Operatoren des Opt-in-Moduls\ZSFTitleTag{11\_math\_advanced}]
  \[
    \begin{aligned}
      &\Ker A,\quad \rang A,\quad \Spur A,\quad \diag(A),\\
      &\spanop(v_1,v_2),\quad \eig(A),\quad \proj_U(v),\\
      &\sig(A),\quad \Adj(A),\quad \Re z,\quad \Im z,\quad
        \asterisk,\quad \oasterisk.
    \end{aligned}
  \]
  % \drawbrace ist der Fall, den \ZSFbraceunder NICHT kann: eine Klammer über
  % ZELLGRENZEN hinweg. Für einen zusammenhängenden Term ist \ZSFbraceunder
  % das richtige Werkzeug (vorgeführt im Kapitel Stilmittel).
  \[
    A=\begin{pmatrix}
      1 & 0 & 4 \\
      \tikzmark[xshift=-\ZSFspaceS]{showcase-brace-left}0 &
      1\tikzmark[xshift=\ZSFspaceS]{showcase-brace-right} & 5
    \end{pmatrix}
  \]
  \drawbrace[brace mirrored]{showcase-brace-left}{showcase-brace-right}
  \annote[yshift=-\ZSFspaceM]{brace-\thebrace}{\ZSFfontNote Pivotspalten}
  \ZSFbraceRoom
\end{statementbox}
